\section{Actual NC separation supported in the original double locus}
\label{sec:hironaka-normal-crossings-double-locus-separation}

Let $H$ be a closed normal-crossings carrier in a complex manifold of
dimension $n$, and write $\nu_H(x)$ for its intrinsic number of local
branches. Its intrinsic double locus is
$D_2(H)=\{x:2\leq\nu_H(x)\}$. This set is closed by the established
upper semicontinuity of the actual branch count. A boundary adapted to
$H$ is an already SNC list whose through members have simultaneous
coordinate equations distinct from the original branches of $H$.
All original strict transforms below are the actual differential
strict transforms.

\begin{lemma}[Pointwise branch count under an actual higher blowup]
\label{hir:general-nc-higher-count-over-image}
\lean{Hironaka.HironakaNormalCrossingsDoubleLocusSeparation.branchMultiplicity_strict_highest_le_image}
\leanok
\uses{hir:general-nc-higher-highest-relative-density,hir:general-nc-intrinsic-branch-multiplicity,hir:general-nc-highest-stratum-global-strict-count-decrease}
Suppose every branch count of $H$ is at most $r\geq2$, and $\pi$ is
an actual blowup of its whole $r$-branch stratum. For every point $z$
of the blowup,
\[
 \nu_{H'}(z)\leq\nu_H(\pi(z)),\qquad H'=\operatorname{Strict}_{\pi}H.
\]
\end{lemma}
\begin{proof}
The established relative density off this higher stratum identifies
the geometric and differential strict transforms. If $\pi(z)$ lies
in the centre, its original count is at least $r$, whereas the actual
strict-carrier count is at most $r-1$ by the genuine monomial blowup
calculation. This proves the inequality there. Otherwise choose a
centred NC presentation of $H$ at $\pi(z)$ with selected index set $S$.
The actual blowup projection is locally biholomorphic off its centre;
pulling this chart back gives a centred chart at $z$ with the same
selected coordinates and literal original membership equation. The
strict transform agrees with the inverse image on that domain.
The intrinsic branch-count dictionary gives both counts as $|S|$,
proving equality off the centre and the asserted inequality everywhere.
\end{proof}

\begin{lemma}[Actual higher towers do not increase branches over their images]
\label{hir:general-nc-higher-tower-count-over-image}
\lean{Hironaka.HironakaNormalCrossingsDoubleLocusSeparation.branchMultiplicity_strict_tower_le_image}
\leanok
\uses{hir:general-nc-higher-count-over-image,hir:general-nc-intrinsic-higher-separation-step}
Let $T$ be any finite actual tower whose centres are genuine whole
attained higher branch-count maxima of the current strict original,
with its literal accumulated adapted boundary. Whenever its terminal
strict original is NC, its branch count at $z$ is at most the original
count at $T(z)$.
\end{lemma}
\begin{proof}
Induct on the actual finite tower. For the identity tower the strict
transform of the original closed carrier is itself, giving equality.
For a last blowup, the recorded higher-centre condition supplies the
current NC carrier, its actual bound $r\geq2$, and the exact whole
centre. The previous lemma bounds the new count by the current count
at the last projection. The induction hypothesis bounds that current
count by the original count at the composite projection. The actual
composition law for differential strict transforms identifies this
two-step strict carrier with the strict carrier of the whole tower.
Transitivity of these two inequalities completes the induction.
\end{proof}

\begin{theorem}[Every actual higher separation centre lies over the double locus]
\label{hir:general-nc-higher-centres-over-double-locus}
\lean{Hironaka.HironakaNormalCrossingsDoubleLocusSeparation.centresOver_doubleLocus_of_higherSeparationCentres,Hironaka.HironakaNormalCrossingsDoubleLocusSeparation.isIsomorphismOver_compl_doubleLocus_of_higherSeparationCentres}
\leanok
\uses{hir:general-nc-higher-tower-count-over-image,hir:general-nc-intrinsic-higher-separation-step}
Every centre of such a tower lies over $D_2(H)$, and its actual
projection is an isomorphism over $D_2(H)^c$.
\end{theorem}
\begin{proof}
Induct through the recorded centres. At a point of a current centre,
the centre equation gives current original count at least its attained
maximum $r\geq2$. The preceding tower inequality bounds this count
above by the original count at its image. The image therefore belongs
to $D_2(H)$; together with the induction hypothesis this proves the
assertion for every centre. The complement of $D_2(H)$ is open.
Restriction of the tower there has every centre empty. The actual
empty-centre restriction theorem then puts the entire exceptional
image in $D_2(H)$. The proved isomorphism theorem for a genuine blowup
sequence outside its exceptional image gives the required isomorphism
on the whole complement, including points on the one-branch locus.
\end{proof}

\begin{lemma}[Actual geometric boundary admissibility]
\label{hir:general-nc-double-locus-geometric-admissibility}
\lean{Hironaka.HironakaNormalCrossingsDoubleLocusSeparation.higherSeparationCentres_isAdmissibleBlowupSequence_doubleLocus}
\leanok
\uses{hir:general-nc-higher-centres-over-double-locus,hir:general-nc-actual-higher-separation-step}
Every such actual higher-separation tower is a genuine admissible
boundary blowup sequence over $D_2(H)$, from the original old boundary
to its literal accumulated strict-old and exceptional list.
\end{lemma}
\begin{proof}
Induct on the tower. The identity is the reflexive admissible
sequence and has exactly the original boundary. At a last blowup,
the recorded higher-centre condition gives an actual whole highest
stratum in simultaneous original and old-boundary coordinates.
Those coordinates prove that the centre has normal crossings with
the current old boundary. Its containment over the original double
locus is the preceding theorem evaluated at that recorded last
centre. Apply the actual admissible-sequence step constructor to
the induction hypothesis and this same blowup, centre, normal-crossing
condition and containment. The resulting boundary is exactly the
strict transform of each existing list member, followed by the
exceptional preimage, which is the literal accumulated boundary.
\end{proof}

\begin{theorem}[A genuine permitted boundary prephase]
\label{hir:general-nc-permitted-double-locus-prephase}
\lean{Hironaka.HironakaNormalCrossingsDoubleLocusSeparation.naturalSeparation_doubleLocus_geometry,Hironaka.HironakaNormalCrossingsDoubleLocusSeparation.naturalSeparation_permitted_geometry}
\leanok
\uses{hir:general-nc-double-locus-geometric-admissibility,hir:general-nc-internally-natural-geometric-separation}
The internally constructed natural separation has at most $n-1$
actual blowups and is a genuine admissible boundary sequence over
$D_2(H)$. Its terminal strict original is smooth and transverse
to its literal accumulated SNC boundary; appending that strict
original is SNC. Every centre lies over $D_2(H)$, and the projection
is an isomorphism on $D_2(H)^c$. More generally, for any closed
permitted set $F$ with $D_2(H)\subseteq F$, the same internally
constructed tower is an admissible boundary sequence with every centre
over $F$ and is an isomorphism
over $F^c$.
\end{theorem}
\begin{proof}
Use the actual natural separation producer driven by the intrinsic
finite range $\{0,\ldots,n\}$ and its internally chosen actual
highest-stratum blowups. Its bounded recursion proves the length,
terminal smoothness, transversality, literal SNC list, and genuine
higher-centre condition. Apply the preceding theorem to that actual
tower and the geometric boundary admissibility lemma. For the final
assertion, each centre's inverse-image
containment in $D_2(H)$ implies its inverse-image containment in $F$.
Restrict to the open complement of the given closed set $F$; all
these actual centres are empty there. The same genuine exceptional
image and isomorphism theorems prove the asserted isomorphism over
$F^c$ for this unchanged tower.
\end{proof}

This boundary prephase is supported precisely within the original
multiple-branch locus. It supplies no marked-ideal cosupport or
ordered-history condition. To use it with a prescribed permitted
bad set, that set must contain the original double locus; this
containment is an explicit geometric hypothesis, not an implicit
change of the original good locus.

\begin{lemma}[One branch is exactly actual hypersurface regularity]
\label{hir:general-nc-one-branch-iff-regular}
\lean{Hironaka.HironakaNormalCrossingsDoubleLocusSeparation.regularAtDim_iff_branchMultiplicity_eq_one}
\leanok
\uses{hir:general-nc-intrinsic-branch-multiplicity,hir:general-nc-count-one-smooth-terminal-pair}
In positive ambient dimension, a point of an actual NC carrier is
regular of dimension $n-1$ if and only if its intrinsic branch count
equals one.
\end{lemma}
\begin{proof}
An actual regular flattening identifies the whole carrier near the
point with the last coordinate hyperplane. It is therefore a genuine
NC presentation with a singleton selected set. The intrinsic branch
dictionary gives count one. Conversely, suppose the intrinsic count
at the point is one. The complement of the closed double locus is an
open neighbourhood of that point. On this actual open submanifold,
the NC carrier is the literal original trace and its branch counts
are the original branch counts, because the inclusion is locally
biholomorphic. All of them are at most one. The point belongs to the
carrier by positivity of its count. The established actual count-one
regularity theorem on this open submanifold supplies dimension
$n-1$ regularity of the trace. The actual open-inclusion regularity
dictionary transfers this regularity to the unchanged original carrier.
\end{proof}

\begin{lemma}[The actual double locus is the singular boundary locus]
\label{hir:general-nc-double-locus-eq-singular-locus}
\lean{Hironaka.HironakaNormalCrossingsDoubleLocusSeparation.doubleLocus_eq_diff_regularLocus}
\leanok
\uses{hir:general-nc-one-branch-iff-regular}
In positive ambient dimension,
\[
 D_2(H)=H\setminus\operatorname{Reg}_{n-1}(H).
\]
\end{lemma}
\begin{proof}
A point of count at least two belongs to the carrier, by count
positivity, and cannot be regular since regularity would give count
one. Conversely a point of the carrier has positive integer count.
If that count is less than two it equals one and is regular by the
preceding equivalence. Every nonregular point of the carrier must
therefore have count at least two. This proves equality pointwise.
\end{proof}

\begin{theorem}[Actual finite NC boundary separation preserves the original smooth locus]
\label{hir:general-nc-original-smooth-locus-snc-prephase}
\lean{Hironaka.HironakaNormalCrossingsDoubleLocusSeparation.originalSeparation_singularLocus_geometry}
\leanok
\uses{hir:general-nc-double-locus-eq-singular-locus,hir:general-nc-permitted-double-locus-prephase,hir:general-nc-generated-prefix-whole-total-support}
Every actual NC carrier in positive dimension internally produces
an actual tower of at most $n-1$ higher-stratum blowups with centres
over its original singular boundary locus. It is a genuine geometric
admissible boundary sequence from the empty old list to its literal
accumulated exceptional list. The terminal strict original is smooth
and transverse to that list; appending the entire strict original
gives a genuine SNC list whose whole union equals the whole original
inverse image. The actual projection is an isomorphism on the entire
complement of the original singular boundary locus.
\end{theorem}
\begin{proof}
Apply the internally generated natural separation to the actual
original vanishing ideal with marking one and empty old boundary.
The preceding prephase theorem proves the finite length, geometric
admissibility, smoothness, transversality, SNC list, centre support
and isomorphism on the complement of the original double locus.
The actual one-branch regularity equivalence identifies that locus
with the original singular boundary locus, so the same tower has
exactly the stated permitted support and good-locus isomorphism.
The established literal original separation support dictionary
identifies the entire preimage of the original carrier with the
union of the unchanged terminal strict original and every accumulated
exceptional member. This is precisely the appended SNC list already
produced by this same original separation, yielding the whole
inverse-image assertion without choosing local branch labels.
\end{proof}
